\section*{Hoofdstuk \ref{ch:b1:observing-the-sky} --- Telescopen, detectoren en fotometrie}

\begin{solution}{exo:b1:observing-the-sky:1}
Een verschil van 5 magnitudes is een factor 100: de ster van magnitude 1.0 zendt
100 keer zoveel flux uit. Rigel en Polaris: $10^{0.4\times(2.02 - 0.13)} =
10^{0.756} = 5.70$; Rigel, met de kleinere magnitude, is helderder.
\end{solution}

\begin{solution}{exo:b1:observing-the-sky:2}
$d = 1/0.19495 = \qty{5.13}{\parsec}$, dat is $5.1295\times206265 =
\qty{1.058e6}{\au}$ en $5.1295\times3.2616 = \qty{16.7}{\lightyear}$.
\end{solution}

\begin{solution}{exo:b1:observing-the-sky:3}
$1.22\times\num{550e-9}/2.4 = \qty{2.80e-7}{rad} = \ang{;;0.058}$, ongeveer zeven keer
fijner dan een \omterm{def:b1:observing-the-sky:seeing}{seeing} van \ang{;;0.40}: in de ruimte bereikt de telescoop zijn
diffractiegrens. Voor \qty{6.5}{m} bij \qty{2.0}{\micro\metre}:
$1.22\times\num{2.0e-6}/6.5 = \qty{3.75e-7}{rad} = \ang{;;0.077}$.
\end{solution}

\begin{solution}{exo:b1:observing-the-sky:4}
$B - V$ stijgt van Rigel ($-0.03$) naar Altair (0.22), Polaris (0.60) en
Betelgeuse (1.85): de heetste is Rigel, daarna Altair, Polaris en Betelgeuse.
De afstand verzwakt beide banden met dezelfde factor, dus $B - V$ verandert er niet
door (zolang geen stof het licht onderweg rood kleurt).
\end{solution}

\begin{solution}{exo:b1:observing-the-sky:5}
$X_1 = \sec\ang{30} = 1.1547$, $X_2 = \sec\ang{60} = 2.0000$.
$k = (6.590 - 6.412)/(2.0000 - 1.1547) = 0.178/0.8453 = \num{0.211}$ per
\omterm{def:b1:observing-the-sky:extinction}{luchtmassa}, en $m_0 = 6.412 - 0.2106\times1.1547 = 6.169$.
\end{solution}

\begin{solution}{exo:b1:observing-the-sky:6}
$\mathrm{SNR} = \sqrt{400} = 20$. De SNR groeit als $\sqrt t$: voor 100 geldt
$t = \qty{10}{s}\times(100/20)^2 = \qty{250}{s}$. Een fout van \num{0.01}
magnitude vereist $\mathrm{SNR} = 1.086/0.01 = 108.6$, dus
$t = 10\times(108.6/20)^2 = \qty{295}{s}$.
\end{solution}

\begin{solution}{exo:b1:observing-the-sky:7}
De afstand is $d = 1/0.00378 = \qty{264.6}{\parsec}$, en de \omterm{def:b1:observing-the-sky:absolute}{absolute magnitude}
$M = 0.13 - 5\log_{10}26.46 = 0.13 - 7.11 = -6.98$. $L/L_\odot = 10^{0.4\times(4.83 + 6.98)} =
10^{4.724} = \num{5.3e4}$ in V.
\end{solution}

\begin{solution}{exo:b1:observing-the-sky:8}
$\pi\times1.2^2 = \qty{4.52}{m^2}$, en $25/4.52 = 5.5$: de Webb-telescoop vangt in een
belichting die 5.5 keer korter is evenveel fotonen op. Een volledige schijf van
\qty{6.5}{m} doorsnede zou $\pi\times3.25^2 = \qty{33.2}{m^2}$ hebben; de werkelijke
spiegel heeft daarvan $25/33.2 = \qty{75}{\percent}$, omdat 18 zeshoeken de rand van die
schijf niet bedekken en de spleten tussen de segmenten en de centrale opening niets
opvangen.
\end{solution}

\begin{solution}{exo:b1:observing-the-sky:9}
(a) $\qty{10}{\parsec} = \qty{3.0857e17}{m}$, en $F = \num{3.828e26}/(4\pi
\times(\num{3.0857e17})^2) = \num{3.828e26}/\num{1.1965e36} =
\qty{3.20e-10}{W.m^{-2}}$, $1361/\num{3.20e-10} = \num{4.3e12}$ keer minder dan bij de aarde.
(b) $5\log_{10}(d/\qty{10}{\parsec}) = -26.74 - 4.83 = -31.57$, dus
$d = \qty{10}{\parsec}\times10^{-6.314} = \qty{4.853e-6}{\parsec}$, dat is
$\num{4.853e-6}\times206265 = \qty{1.001}{\au}$.
\end{solution}

\begin{solution}{exo:b1:observing-the-sky:10}
Boven de atmosfeer is de flux $F_0$; daaronder $F = F_0e^{-\tau_0X}$. Dan geldt
$m - m_0 = -2.5\log_{10}(F/F_0) = 2.5\log_{10}e\;\tau_0X$, dat is $kX$ met
$k = 2.5\log_{10}e\;\tau_0 = 1.0857\,\tau_0$. Voor $k = 0.160$:
$\tau_0 = 0.160/1.0857 = 0.147$. In het \omterm{def:b1:spherical-astronomy:celestial-sphere}{zenit} komt $e^{-0.147} = 0.863$ van het licht aan,
\qty{86}{\percent}; bij $X = 2.31$ is dat $e^{-0.147\times2.31} =
e^{-0.340} = 0.712$, \qty{71}{\percent}. Bij $z = \ang{85}$ geldt $\sec z =
11.5$; maar het vlakparallelle model maakt het pad oneindig aan de horizon, terwijl het pad
door een gekromde atmosfeer eindig blijft, en ook breking buigt de straal: nabij de horizon
overschat de formule de \omterm{def:b1:observing-the-sky:extinction}{luchtmassa}, en de fouten groeien snel.
\end{solution}

\begin{solution}{exo:b1:observing-the-sky:11}
(a) $\mathrm{SNR} = Rt/\sqrt{(R + B)t}$, dus $t = \mathrm{SNR}^2(R + B)/R^2
= 100\times2050/2500 = \qty{82}{s}$. (b) Halvering van de \omterm{def:b1:observing-the-sky:seeing}{seeing} deelt het oppervlak van het beeld,
en dus $B$, door 4: $B = 500$ en $t = 100\times550/2500 =
\qty{22}{s}$. (c) Voor $B \gg R$ geldt
$t \approx \mathrm{SNR}^2B/R^2$. De bronsnelheid $R$ groeit als $D^2$; de hemelsnelheid $B$ groeit
als het opvangoppervlak $D^2$ maal het hemeloppervlak van het beeld, $\theta^2$. Dus
$t \propto
\theta^2D^2/D^4 = \theta^2/D^2$: de \omterm{def:b1:observing-the-sky:seeing}{seeing} halveren is even veel waard als de
diameter verdubbelen.
\end{solution}

\begin{solution}{exo:b1:observing-the-sky:12}
(a) $N = 1200\times100 = \num{120000}$ lijnen; in de eerste orde is het scheidend vermogen
\num{1.2e5}, en $\Delta\lambda = 550/(\num{1.2e5}) =
\qty{4.6e-3}{nm}$. (b) $v = c\,\Delta\lambda/\lambda = c/(\num{1.2e5}) =
\qty{2.5}{km.s^{-1}}$. (c) $\Delta\lambda = 550\times29.78/\num{299792} =
\qty{0.055}{nm}$, ongeveer 12 resolutie-elementen. De waarnemer moet voor de datum en het
tijdstip van de waarneming de component van de baansnelheid van de aarde naar de ster toe
berekenen en aftrekken, om de snelheid van de ster ten opzichte van de zon te krijgen.
\end{solution}

\begin{solution}{pb:b1:observing-the-sky:1}
\textbf{1.} $d = 1/0.37921 = \qty{2.637}{\parsec}$.

\textbf{2.} $2.637\times206265 = \qty{5.44e5}{\au}$;
$2.637\times\num{3.0857e16} = \qty{8.14e16}{m}$.

\textbf{3.} $2.637\times3.2616 = \qty{8.60}{\lightyear}$: het licht verliet Sirius 8.6 jaar geleden.

\textbf{4.} De richting beschrijft $2p = \ang{;;0.758}$, ongeveer twee keer de
\omterm{def:b1:observing-the-sky:seeing}{seeing}. De uitslag wordt gemeten vanaf de plaats van het middelpunt van het beeld, niet
vanaf de rand: gemiddeld over veel fotonen wordt het middelpunt tot op een kleine fractie
van de breedte van het beeld bepaald, des te nauwkeuriger naarmate er meer fotonen zijn geteld.

\textbf{5.} $\Delta d/d = \Delta p/p = 0.001/0.37921 = \qty{0.26}{\percent}$,
dat is $8.60\times0.0026 = \qty{0.023}{\lightyear}$.

\textbf{6.} $10^{0.4\times(0.03 + 1.46)} = 10^{0.596} = 3.94$.

\textbf{7.} $10^{0.4\times(26.74 - 1.46)} = 10^{10.112} = \num{1.29e10}$.

\textbf{8.} $10^{0.4\times(8.44 + 1.46)} = 10^{3.96} = \num{9120}$.

\textbf{9.} $z = \ang{90} - \ang{82.09} = \ang{7.91}$, $X = \sec\ang{7.91} =
1.0096$; extinctie $0.15\times1.0096 = 0.151$ magnitude; gemeten
$m = -1.46 + 0.151 = -1.31$.

\textbf{10.} $z = \ang{60}$, $X = 2.000$; extinctie 0.30 magnitude,
$m = -1.16$. De doorgelaten fractie is $10^{-0.4\times0.30} = 0.759$:
\qty{24}{\percent} van het licht gaat verloren.

\textbf{11.} Bij de opkomst is $z = \ang{90}$ en is $\sec z$ oneindig, terwijl het
werkelijke pad door de gekromde atmosfeer lang maar eindig is; ook breking buigt de straal.
Nabij de horizon is de extinctie groot en verandert zij snel, en de vlakparallelle correctie
is onjuist.

\textbf{12.} $\num{3.0e4}\times10^{0.4\times(10.0 - 8.44)} = \num{3.0e4}
\times10^{0.624} = \num{3.0e4}\times4.21 = \num{1.26e5}$ fotoelektronen per
seconde.

\textbf{13.} Een SNR van 1000 vereist $N = 1000^2 = \num{1.0e6}$
fotoelektronen: $t = \num{1.0e6}/\num{1.26e5} = \qty{7.9}{s}$. De fout op de magnitude
is $1.086/1000 = \num{0.0011}$.

\textbf{14.} $\num{3.0e4}\times10^{0.4\times(10.0 + 1.46)} = \num{3.0e4}
\times10^{4.584} = \num{1.15e9}$ fotoelektronen per seconde.

\textbf{15.} $0.20\times\num{1.15e9} = \num{2.30e8}$ per seconde in de helderste pixel, die
in $\num{1.0e5}/\num{2.30e8} =
\qty{4.3e-4}{s}$ vol is, minder dan een halve milliseconde.

\textbf{16.} De vleugels van Sirius geven $\num{1.15e9}\times10^{-4} =
\num{1.15e5}$ per seconde; Sirius~B geeft $0.20\times\num{1.26e5} =
\num{2.52e4}$ per seconde. De gloed van Sirius is 4.6 keer helderder dan de piek van Sirius~B.

\textbf{17.} De pixel ontvangt $\num{1.15e5} + \num{2.52e4} = \num{1.40e5}$
per seconde en is in $\num{1.0e5}/\num{1.40e5} = \qty{0.71}{s}$ vol. In
\qty{0.5}{s}: signaal $\num{2.52e4}\times0.5 = \num{1.26e4}$, totaal
$\num{7.01e4}$, ruis $\sqrt{\num{7.01e4}} = 265$, dus $\mathrm{SNR} =
\num{1.26e4}/265 = 48$. Sirius~B is detecteerbaar voor zover het fotonruis betreft; wat haar
moeilijk maakt, is dat de gloed bekend moet zijn en moet worden afgetrokken tot veel beter dan
een vijfde van haar waarde, terwijl de kern van Sirius in minder dan een milliseconde verzadigt.

\textbf{18.} $m - M = 5\log_{10}(2.637/10) = 5\times(-0.579) = -2.89$.

\textbf{19.} $M = -1.46 + 2.89 = 1.43$.

\textbf{20.} $d = 1/0.13023 = \qty{7.679}{\parsec}$;
$m - M = 5\log_{10}0.7679 = -0.57$, dus $M = 0.03 + 0.57 = 0.60$.

\textbf{21.} Vega, met de kleinere \omterm{def:b1:observing-the-sky:absolute}{absolute magnitude}, is helderder met de factor
$10^{0.4\times(1.43 - 0.60)} = 10^{0.332} = 2.15$. Sirius lijkt alleen helderder omdat zij
dichterbij staat.

\textbf{22.} $m = M_\odot + (m - M) = 4.83 - 2.89 = 1.94$: de zon zou een tamelijk heldere
ster met het blote oog zijn.

\textbf{23.} $5\log_{10}(d/\qty{10}{\parsec}) = 6.0 - 1.43 = 4.57$, dus
$d = 10\times10^{0.914} = \qty{82.0}{\parsec}$, ongeveer \qty{267}{\lightyear}.

\textbf{24.} $L/L_\odot = 10^{0.4\times(4.83 - 1.43)} = 10^{1.36} = 22.9$:
in de V-band overstraalt Sirius \textbf{ongeveer 23 zonnen}.
\end{solution}
